\documentclass[instructions.tex]{subfiles}
\begin{document}

\chapter{Du péché véniel.}


Il n'y a qu'une chose en ce monde, dit saint Jean Chrysostome, qu'on doive craindre, c'est le péché, quel qu'il soit; il est l'unique mal : tous les autres ne sont que des bagatelles\footnote{Unica duntaxat res gravis et perlimiscenda, peccatum, scilicet; reliqua verò omnia mera fabula. S. Joan. Chrys.}. Le péché véniel, à la vérité, n'est qu'un petit péché; mais un petit péché ne laisse pas d'être un grand mal, soit par rapport à Dieu, soit par rapport à nous.

\primo Tout péché offense Dieu, est opposé à sa volonté, défendu par sa loi, et blesse sa sainteté : il est donc comme le mal de Dieu. Or, le moindre mal de Dieu est un plus grand mal que tous les maux des créatures. Le moindre péché a affligé Jésus-Christ, et lui a fait verser des larmes. Devons-nous regarder comme un rien ce qui a causé de l'affliction à un Dieu? Si l'on se fait une peine de causer du chagrin à un ami, pourquoi sommes-nous si malheureux que d'affliger si souvent un Dieu qui nous aime plus tendrement que tous nos amis ne peuvent nous aimer?

Nous comprendrions ce que c'est qu'un péché véniel, si nous avions les sentiments des saints. Oh! qu'ils se croient heureux dans le ciel de ne plus offenser Dieu, et plus heureux encore d'être assurés que jamais ils ne pourront l'offenser! Ils connaissent si parfaitement l'injure qu'un léger péché fait à Dieu, qu'ils aimeraient mieux souffrir les tourments de l'enfer, que de se rendre coupables, s'il était possible, d'un seul péché véniel.

Combien de fois Dieu lui-même a-t-il puni exemplairement des fautes que nous regardons comme légères! La femme de Loth, pour un regard curieux, frappée de mort. Moïse, l'ami de Dieu, pour une légère défiance, exclu de la terre promise à ses pères. Le royaume de Juda désolé par une peste cruelle, pour une vanité de David. Oza frappé de mort, pour avoir porté sa main sur l'arche sainte. Ézéchias, ce saint roi, puni de Dieu jusqu'à la troisième génération, pour une complaisance indiscrète envers les ambassadeurs étrangers. Quarante-deux enfans punis de mort, pour une raillerie contre un prophète. Un Israélite condamné à mort par l'ordre de Dieu, pour avoir recueilli un peu de bois un jour de fête. Le roi Ozias puni de Dieu, pour avoir touché à l'encensoir.

Mais en l'autre vie, la peine qui doit expier le péché véniel en purgatoire est si incompréhensible, que tous les tourmens d'ici-bas n'ont rien qui lui soit comparable. Les fautes qui nous paraissent légères sont donc un grand mal par rapport à Dieu, puisqu'il les punit avec tant de rigueur.

\secundo Ces fautes sont encore un grand mal par rapport à nous. Il est vrai que le péché véniel ne nous sépare pas de Dieu; mais il indispose Dieu, et refroidit l'amour qu'il a pour nous. Il ne nous fait pas perdre la grâce sanctifiante, mais il en ternit l'éclat : il empêche, il affaiblit ou rend inutiles les secours que Dieu nous destinait pour persévérer. Le péché véniel n'exclut pas du ciel, mais il nous en éloigne. Il ne damne pas, mais il dispose peu à peu à la damnation : il nous approche du précipice; et quand on est sur le bord, que faut-il, hélas! pour y tomber? Il ne donne pas le coup de la mort à l'âme, mais il est comme une langueur qui conduit peu à peu à la mort. Oh! si l'on réfléchissait à ces vérités, avec quelles précautions éviterait-on les moindres fautes!

Disons ici pour la consolation de ceux qui cherchent Dieu sincèrement, que les justes mêmes ne sont pas exempts de ces fautes légères; et que celles qu'on commet par surprise, par fragilité, contribuent à notre perfection, quand on s'en humilie et qu'on en gémit. Mais on doit dire aux âmes lâches, que les péchés véniels que l'on commet habituellement de propos délibéré, et qu'on ne compte pour rien, mettent une âme dans l'état le plus dangereux.

Mépriser les petites fautes, c'est, selon l'oracle du Saint-Esprit, s'exposer à de grandes chutes, et s'exposer à périr\footnote{Qui spernit modica, paulatim decidet. Eccl. 19.}. Un peu de vanité, de curiosité; un peu trop de liberté dans ses paroles, d'attache à ses biens, à ses aises; un peu d'antipathie, de jalousie, ou d'affection trop naturelle; un peu trop de complaisance et de respect humain; certains désirs de paraître, certaine négligence à veiller sur son humeur, etc., tout cela, quand il est volontaire et habituel, va plus loin qu'on ne pense. Combien de personnes, hélas! ont commencé leur réprobation en méprisant des fautes légères qui ont commencé leur égarement!

Il importe peu, quand on perd la vie, que ce soit par un coup de foudre ou par une fièvre lente qui affaiblit peu à peu le principe de la vie. Il importe peu, quand on a fait naufrage, que ce soit par un coup de tempête qui engloutit le vaisseau, ou qu'il coule à fond en prenant peu à peu l'eau de tous côtés. De même il importe peu à l'ennemi du salut, quand il veut perdre une âme, de la faire tomber tout-à-coup dans de grands désordres, ou de la conduire peu à peu au précipice.

Nous sommes d'autant plus téméraires de négliger les fautes légères, que nous prenons quelquefois pour véniel ce qui est mortel, et d'autant plus coupables, que nous pouvons, sans beaucoup de peine, nous préserver d'y tomber si fréquemment. Quelle honte de nous laisser vaincre si souvent par notre humeur et par le démon, puisqu'il est si aisé de remporter la victoire! Nous pouvons à chaque moment nous approcher du ciel, et nous nous en éloignons presque à chaque pas. Quelle lâcheté de blesser si souvent son âme, et de déplaire à Dieu par tant de péchés qu'il est facile d'éviter!


\section{Résolutions}

\primo Je ne regarderai plus le péché véniel comme un mal indifférent, puisqu'il offense Dieu, et qu'il fait de si grands torts à l'âme qui s'y livre sans retenue et sans crainte.

\secundo Je m'interdis donc dès à présent ces petits mensonges, ces railleries légères, ces animosités secrètes qui ne vont pas à la perte de la charité; ces négligences dans le service de Dieu et les devoirs de mon état, qui me semblaient peu importantes; ces paroles indiscrètes sur tant de choses où ma conscience ne se croyait pas grièvement blessée; ces gourmandises fréquentes, etc.

Mon Dieu, inspirez-moi un grand éloignement pour tout ce qui vous déplaît et qui est pernicieux à mon âme.

\end{document}
